% from main.tex
%%%%%%%%%%  scan idx 617  |  folio 604  %%%%%%%%%%
soit exacte).

Posons \(A=\mathcal O_Y\), \(B_0=f_*(\mathcal O_X)\) et définissons par récurrence \(B_i=\Ker\bigl(B_{i-1}\rightrightarrows B_{i-1}\otimes_A B_{i-1}\bigr)\). On vérifie aussitôt que \(B_{i-1}\otimes_A B_{i-1}\xrightarrow{\sim}B_{i-1}\otimes_{A_i}B_{i-1}\). Reste à prouver que pour \(i\) assez grand, on a \(B_i=A\). En effet considérons la suite décroissante de \(B\)-modules \(M_i=B_i/A\) et soit \(S_i=\Supp(M_i)\). La suite décroissante de fermés \(S_i\) de \(Y\) est stationnaire, soit \(S\) la valeur stationnaire. Montrons que \(S=\varnothing\). Sinon soit \(s\) un point maximal de \(S\). Quitte à remplacer \(Y\) par \(\Spec\mathcal O_{Y,s}\), on peut supposer \(Y\) local de point fermé \(s\). Vu le choix de \(s\), \(M_i\) est artinien pour \(i\) assez grand, donc la suite \((M_i)\) est stationnaire. Il existe donc \(i\) tel que \(B_i=B_{i+1}\). Par suite \(\Spec(A_i)\to Y\) est un monomorphisme, donc une immersion fermée (Nakayama); comme \(B\subset A_i\), on a \(B=A_i\), et \(s\notin S\), absurde.

\section*{3. Démonstration de I : le dévissage de Oort}

\textbf{Proposition 3.1.} (Oort, \emph{Bull. soc. math. France}, t. 90-1962). Soient \(S\) un schéma localement noethérien, \(f:X\to S\) un morphisme propre, \(X'\subset X\) un sous-schéma fermé. On suppose que~:
\begin{enumerate}[label=\roman*)]
\item \(X'\) est défini par un idéal \(I\) contenu dans le nilradical \(N\) de \(X\) et \(N.I=0\).
\item \(X'\) est plat sur \(S\) et cohomologiquement plat sur \(S\) en dimension \(0\) (EGA III 7.8.1).
\item \(I\) est plat sur \(S\) et cohomologiquement plat sur \(S\) en dimension \(0\) et \(1\) de sorte que (d'après EGA III 7.4.6) il existe un \(\mathcal O_S\)-module cohérent \(Q\) et un isomorphisme fonctoriel en le \(\mathcal O_S\)-module cohérent \(M\),
